% Introduction (draft; results-independent paragraphs only. The "here we show"
% paragraph is written after Gate G3.)

Water levels in inland seas and lakes will fall, often dramatically, through the twenty-first century as the climate warms\citep{prange2020}. Where a closed sea is shallow and its floor nearly flat, even a modest fall exposes large areas. Where that newly exposed seabed lies decides what is lost: harbours and navigation channels, coastal wetlands, and the breeding and feeding grounds of the sea's fauna.

The Caspian Sea is now living through this. By the end of 2024 its level had fallen to about \metric{lit.ginzburg2025.level_end2024_bhs_m}~m in the Baltic Height System, below the 1977 minimum of \metric{lit.ginzburg2025.min1977_bhs_m}~m, after falling by \metric{lit.ginzburg2025.fall_rate_2020_2024_cm_per_yr}~cm per year over 2020--2024. That level is satellite altimetry lowered by about \metric{lit.ginzburg2025.altimetry_offset_m}~m, the amount by which altimetry exceeds instrumental measurements\citep{ginzburg2025}. Projections of the fall still to come differ widely. One CMIP6-based analysis puts the total fall over 2020--2100 at about \metric{lit.samant2023.ssp245_m}~m under SSP2-4.5 (inter-model range \metric{lit.samant2023.ssp245_lo_m}--\metric{lit.samant2023.ssp245_hi_m}~m) and about \metric{lit.samant2023.ssp585_m}~m under SSP5-8.5 (\metric{lit.samant2023.ssp585_lo_m}--\metric{lit.samant2023.ssp585_hi_m}~m). Each total adds two parts: the observed 1996--2021 decline carried forward at a constant water budget, about \metric{lit.samant2023.trend_2100_m}~m by 2100, and a further \metric{lit.samant2023.pe_ssp245_m}~m (SSP2-4.5) or \metric{lit.samant2023.pe_ssp585_m}~m (SSP5-8.5) that a best-fit subset of the models gives from the trend in precipitation minus evaporation\citep{samant2023}. A study of \metric{lit.hoseini2025.n_gcms} CMIP6 models finds much smaller falls, with median changes from 2021 to 2100 of \metric{lit.hoseini2025.medium_median_m}~m for a medium-emission scenario and \metric{lit.hoseini2025.high_median_m}~m for a high-emission one, with 25th--75th percentile ranges of \metric{lit.hoseini2025.medium_q25_m} to \metric{lit.hoseini2025.medium_q75_m}~m and \metric{lit.hoseini2025.high_q25_m} to \metric{lit.hoseini2025.high_q75_m}~m\citep{hoseini2025}.

Much of the coming emergence will happen in the North Caspian. This shelf has a mean depth of \metric{lit.akbari2020.nc_depth_lo_m}--\metric{lit.akbari2020.nc_depth_hi_m}~m and holds less than \metric{lit.akbari2020.nc_volume_share_pct}\% of the sea's volume\citep{akbari2020}. Landsat water masks from 1977--2018 show the flooded part of the band between its 1977 and 1995 shorelines changing by about \metric{lit.akbari2020.nc_area_per_m_km2}~\kmsq{} per metre of level, but that relation was fitted between levels of \metric{lit.akbari2020.fit_swl_lo_m} and \metric{lit.akbari2020.fit_swl_hi_m}~m, and its authors caution against extrapolating it\citep{akbari2020}. The sea has now fallen past the lower end of that range, the 1977 minimum. Below it, projections of how much more of the shelf will emerge, and of what it will take with it, need the shape of the seabed, and recent assessments take it from global relief grids. One, based on GEBCO\_2022, found that a decline of just \metric{lit.court2025.decline_lo_m}--\metric{lit.court2025.decline_hi_m}~m would reduce existing marine protected area coverage by up to \metric{lit.court2025.mpa_loss_max_pct}\%\citep{court2025}. Another took the sea's topography from GEBCO and concluded that a fall to \metric{lit.duan2025.level_2100_m}~m, which it expects by 2100 if the 1995--2023 trend continues, would remove the northern Caspian\citep{duan2025}.

These grids are least certain exactly where the projections lean on them hardest. Globally, only about \metric{lit.weatherall2015.control_share_pct}\% of the non-land cells of GEBCO\_2014 were based on bathymetric control data, that is, soundings or pre-prepared grids that may contain interpolated values\citep{weatherall2015}. In the current GEBCO grid, the North Caspian seabed comes from a pre-generated grid of mixed source data, a type that GEBCO groups with unknown sources, and its vertical reference there is not documented\citep{gebco2026}. On a shelf this flat, a small vertical error moves the predicted shoreline a long way horizontally. Yet the grids have not, to our knowledge, been tested quantitatively, coast by coast, against the seabed that has actually emerged. One study noted that the band flooded as the level rose by about \metric{lit.akbari2020.rise_1977_1995_m}~m between 1977 and 1995 matched depths of less than \metric{lit.akbari2020.vulnerable_depth_m}~m in its ETOPO1 depth map\citep{akbari2020}. Another found areas from a chart-based model within \metric{lit.bocharov2026.max_area_diff_pct}\% of satellite-observed areas for the whole sea\citep{bocharov2026}. The first is a qualitative check; the second is a whole-sea total, in which errors of opposite sign along different coasts can cancel.

A falling closed sea surveys its own floor. Each cloud-free satellite image records where the water's edge lay, and satellite altimetry records how high the water stood that day, so every dated waterline is a contour of the seabed. Waterlines paired with water levels underpin intertidal topography mapped from tidal stages\citep{sagar2017,bishoptaylor2019}. In a closed sea the water level itself sweeps the shelf over years rather than hours. Since the altimetry record began on \metric{levels.dahiti.first}, the Caspian has risen to a maximum on \metric{levels.dahiti.smoothed.date_of_max} and then fallen by about \metric{levels.dahiti.smoothed.range_m}~m to a low on \metric{levels.dahiti.smoothed.date_of_min} (dates and range from a \metric{levels.dahiti.smoothed.window_records}-record running median of the series), crossing the upper part of that range twice.
